\begin{examplebox}
\textbf{\textcolor{CExample}{例 1（基础速算：已知切点求切线）}}

\textbf{题目：} 椭圆 $\dfrac{x^2}{16}+\dfrac{y^2}{9}=1$ 上有点
$P\left(2,\dfrac{3\sqrt3}{2}\right)$。
求 $P$ 处切线的斜率及方程。

\textbf{解：} 先验点：$\dfrac{2^2}{16}+\dfrac{(3\sqrt3/2)^2}{9}
=\dfrac14+\dfrac34=1$，确实在椭圆上。
由 $a^2=16,b^2=9$ 得
\[
 e^2-1=-\frac9{16},\qquad
 k_{OP}=\frac{3\sqrt3/2}{2}=\frac{3\sqrt3}{4}.
\]
直接用乘积公式：
\[
 k_l=\frac{-9/16}{3\sqrt3/4}=-\frac{\sqrt3}{4}.
\]
再用切点式求直线方程：
\[
 \frac{2x}{16}+\frac{(3\sqrt3/2)y}{9}=1
 \quad\Longrightarrow\quad
 \boxed{3x+4\sqrt3\,y=24}.
\]
\textbf{考场技巧：} 只问切线方程时，直接写切点式更快；
只问斜率时，用 $k_{OP}k_l=-9/16$ 更快。

\medskip
\textbf{\textcolor{CExample}{例 2（逆向应用：已知离心率与切线斜率）}}

\textbf{题目：} 椭圆离心率 $e=\dfrac12$。椭圆上一非轴点
$P$ 处切线斜率为 $-\dfrac12$，求 $OP$ 的斜率。

\textbf{解：} $e^2-1=\dfrac14-1=-\dfrac34$，故
\[
 k_{OP}=\frac{e^2-1}{k_l}
       =\frac{-3/4}{-1/2}=\boxed{\frac32}.
\]
\textbf{检查：} 两斜率一正一负，与第一、三象限示意相符；
切点也可能在第三象限，单凭斜率不能确定是哪一个切点。

\medskip
\textbf{\textcolor{CExample}{例 3（求离心率：识别乘积即秒解）}}

\textbf{题目：} 某标准横轴椭圆上一非轴点 $P$ 满足
$k_{OP}=2$，其切线斜率为 $k_l=-\dfrac38$。
求离心率及短长半轴之比。

\textbf{解：} 相乘得
\[
 k_{OP}k_l=2\left(-\frac38\right)=-\frac34=e^2-1,
\]
\[
 e^2=\frac14,\quad \boxed{e=\frac12},\qquad
 \frac{b^2}{a^2}=\frac34,\quad
 \boxed{\frac ba=\frac{\sqrt3}{2}}.
\]
\textbf{提醒：} $e$ 是非负数，不能取 $-1/2$。

\medskip
\textbf{\textcolor{CExample}{例 4（综合变式：直线 $OP$ 已知，反求切点）}}

\textbf{题目：} 椭圆 $\dfrac{x^2}{16}+\dfrac{y^2}{9}=1$，
在第一象限的切点 $P$ 满足 $OP$ 的斜率为 $1$。
求切点坐标和切线方程。

\textbf{解：} $k_{OP}=1$ 表示 $P$ 在 $y=x$ 上，设 $P=(t,t)$，$t>0$。
代入椭圆：
\[
 \frac{t^2}{16}+\frac{t^2}{9}=1
 \quad\Longrightarrow\quad
 \frac{25t^2}{144}=1
 \quad\Longrightarrow\quad t=\frac{12}{5}.
\]
所以 $P=(12/5,12/5)$。
斜率不必再求导：
$k_l=(e^2-1)/k_{OP}=-9/16$。
切点式给出
\[
 \frac{(12/5)x}{16}+\frac{(12/5)y}{9}=1
 \quad\Longrightarrow\quad \boxed{9x+16y=60}.
\]
\textbf{迁移：} 若给 $k_{OP}=r$，可设 $y=rx$ 与椭圆联立，
求出切点，再用已知切线公式。

\medskip
\textbf{\textcolor{CExample}{例 5（防误用：坐标轴顶点）}}

\textbf{题目：} 椭圆 $\dfrac{x^2}{16}+\dfrac{y^2}{9}=1$ 在
$A(4,0)$ 和 $B(0,3)$ 处的切线分别是什么？

\textbf{解：} 直接代入\emph{切点式}：
\[
 A:\ \frac{4x}{16}=1\Longrightarrow\boxed{x=4};\qquad
 B:\ \frac{3y}{9}=1\Longrightarrow\boxed{y=3}.
\]
在 $A$ 处切线斜率不存在，在 $B$ 处 $OB$ 斜率不存在，
所以这两点\textbf{不能直接代入斜率乘积公式}。

\medskip
\textbf{\textcolor{CExample}{高价值变式（自检）}}
若某标准横轴椭圆上一非轴点满足
$k_{OP}k_l=-\dfrac49$，不求切点，能否直接求离心率？
\textbf{答案：} $e^2=1-4/9=5/9$，故
$\boxed{e=\sqrt5/3}$。\textbf{识别信号：} 两条斜率的\emph{乘积}给出时，
不必引入 $x_0,y_0$。
\end{examplebox}
